\section{Kognitive Architektur (Kap. 02) -- Vertiefungskapitel}
\label{sec:kognitive-architektur}

\subsection{Ist-Architektur: Perception, Memory, Planning und Action}

Als Beobachtungsrahmen wird zwischen Perception, Memory, Planning und Action
unterschieden. Diese Zuordnung ist eine konzeptionelle Analyse des Codes, keine
Behauptung, dass vier eigenständige Laufzeitmodule existieren. CoALA präzisiert
für Sprachagenten insbesondere explizites Arbeits- und Langzeitgedächtnis,
interne und externe Aktionen sowie einen wiederholten Entscheidungszyklus
\parencite{sumersEtAl2024}. StudyMummy realisiert Teile davon als
FastAPI-Endpunkt, Datenbankservices und drei LLM-Rollen.

Der konkrete Kontrollfluss beginnt in \texttt{chat()}: Eingabe filtern,
Eigentum an Aufgabe und Dokument prüfen, Sitzungszustand laden und Hilfestufe
aktualisieren, Nutzerturn speichern, RAG-Kontext abrufen und den typisierten
\texttt{AgentContext} bilden. Der Orchestrator ruft danach Planner, Tutor und
optional Reviewer auf. Erst außerhalb des Orchestrators speichert der Endpunkt
die finale Antwort und ergänzt den Trace-Schritt \enquote{remember}. Dieser
letzte Punkt ist relevant: Persistenz ist vorhanden, aber nicht Teil einer
lernfähigen Rückkopplung des Orchestrators.

\subsubsection{Perception}

Die Rohwahrnehmung besteht aus freier Nutzereingabe und Datenbankzustand.
\texttt{filter\_user\_input()} normalisiert Unicode, entfernt unsichtbare
Steuerzeichen, reduziert Leerraum, begrenzt auf 4.000 Zeichen und blockiert
einige reguläre Ausdrücke für direkte Prompt-Injection. Anschließend verknüpft
\texttt{\_resolve\_study\_context()} die Nachricht nur mit Aufgaben und
Dokumenten des authentifizierten Nutzers. Die Hilfestufe wird anhand weniger
Schlüsselwörter deterministisch auf 1 bis 4 angehoben. Für RAG wird die Frage
mit dem vertrauenswürdigen Aufgabenkontext kombiniert, der Abruf filtert nach
\texttt{user\_id} und optional \texttt{document\_id}.

Die resultierende Repräsentation ist zweckmäßig, aber selektiv. Der Planner
erhält weder Dokumentinhalt noch Dialoghistorie, sondern nur Nachricht,
Hilfestufe, Aufgabe und ein boolesches Kontextmerkmal. Das reduziert
Injection-Risiken, verhindert aber dokumentabhängige Planung. Umgekehrt kann
\texttt{extra\_context} bis zu 12.000 Zeichen aus dem Frontend enthalten und
wird nicht durch den Eingabefilter geprüft. RAG-Text wird im Tutor-Prompt zwar
als nicht vertrauenswürdig markiert, eine strukturelle Trennung von Daten und
Anweisungen findet jedoch nicht statt. Fehlgeschlagener Retrieval-Zugriff führt
laut \texttt{rag\_service.py} still zu leerem Kontext und ist für Planner und
Nutzer nicht unterscheidbar von einem legitimen Treffer ohne Ergebnis.

\subsubsection{Memory}

Das Arbeitsgedächtnis ist explizit in \texttt{Session} und
\texttt{WorkingMemory} modelliert: aktive Aufgabe, Hilfestufe und kompletter
Dialog werden vor jedem Turn geladen. In CoALA ist ein solches
Arbeitsgedächtnis eine über LLM-Aufrufe persistente Datenstruktur und nicht nur
das Kontextfenster \parencite{sumersEtAl2024}. StudyMummy erfüllt dies teilweise,
kopiert jedoch die gesamte Historie ohne Tokenbudget, Zusammenfassung oder
Vergessen in jeden Tutoraufruf. Lange Sitzungen wachsen daher unbeschränkt bis
zum Providerlimit.

Episodisch persistieren \texttt{ChatLog}-Datensätze Rolle, Inhalt, Zeitpunkt
und die finale Aktion. Pläne, Reviewer-Feedback, Werkzeugargumente,
Werkzeugresultate und Erfolg oder Fehlschlag werden nicht gespeichert und
können später nicht gezielt abgerufen werden. Semantisches Gedächtnis besteht
aus hochgeladenen Dokumenten: Der Analyzer zerlegt Volltext mit
\texttt{textwrap.wrap} in ungefähr 1.000 Zeichen lange Chunks, erzeugt
Embeddings und speichert 1.536-dimensionale Vektoren in PostgreSQL. Online
werden die drei Chunks mit kleinster Kosinusdistanz geladen. Eine
Ähnlichkeitsschwelle, ein Reranker, überlappende oder strukturerhaltende Chunks
und Quellenzitate fehlen. Lernprofile und Confidence Scores bilden zwar einen
weiteren Langzeitspeicher, bleiben im analysierten Chatpfad aber ungenutzt.

\subsubsection{Planning}

Der Planning Agent erzeugt mit Temperatur 0 ein \texttt{AgentPlan}-Objekt aus
Intent, genau einer Aktion, Ziel, kurzer Entscheidungsbasis, Werkzeugnamen und
Erfolgskriterien. Pydantic validiert das JSON. Bei Provider- oder
Strukturfehlern greift ein deterministischer, schlüsselwortbasierter Fallback.
Danach schneidet \texttt{SAFE\_TOOL\_POLICY} die vorgeschlagenen Rechte auf ein
aktionsabhängiges Set zurück. Diese Kombination ist nachvollziehbarer als ein
freier Prompt und implementiert eine begrenzte deliberative Auswahl.

Sie ist jedoch weder hierarchische Zielzerlegung noch Closed-Loop-Planung. Der
Planner kennt Werkzeugbeobachtungen nicht, bewertet keine Alternativen und wird
nach Toolfehler oder Reviewer-Ablehnung nicht erneut aufgerufen. Der Tutor darf
innerhalb einer Aktion bis zu fünf Function-Calling-Runden ausführen, diese
lokale Schleife ähnelt ReAct, während der übergeordnete Plan unverändert bleibt.
ReAct begründet gerade den Nutzen, Reasoning und externe Beobachtung
abwechselnd zur Planaktualisierung zu verwenden \parencite{yaoEtAl2023react}.
StudyMummy nutzt die Beobachtung nur im Tutor-Dialog, nicht im typisierten
Planungszustand.

\subsubsection{Action}

Action umfasst Textausgabe und echte Function Calls. Der LLM-Adapter bietet nur
die freigegebenen JSON-Schemas an, blockiert unbekannte Namen erneut, parst
Argumente und gibt Werkzeugresultate als \texttt{tool}-Nachricht an den Tutor
zurück. JSON-Fehler und Python-Ausnahmen werden in eine Fehlerbeobachtung
umgewandelt. Praktisch erreichbar sind Antwortbewertung, Vergabe virtueller
Coins/Erfahrung und ein Kalendereintrag. Die Belohnung bindet die Identität über
einen Request-Kontext, begrenzt den Betrag, prüft Aufgabeneigentum unter
Datenbanksperre und verhindert Doppelbelohnung.

Der Reviewer sieht anschließend nur Plan und fertigen Entwurf. Er kann Text
ersetzen, aber einen bereits angelegten Termin oder eine Belohnung nicht
zurücknehmen. Bei ungültigem Reviewer-JSON wird der Entwurf sogar freigegeben.
Außerdem verlangt der Tutor-Prompt ein Lernprofil-Update, obwohl dieses Tool
laut Policy nie im Scope liegt. Diese Diskrepanz zwischen prozeduralem Wissen
(Prompt) und ausführbarem Aktionsraum ist ein konkreter Architekturfehler.

\subsection{Kritische Analyse und Grenzen}

Die Schichten sind grundsätzlich erkennbar, aber ihre Rückkopplung ist schwach.
Erstens ist die persistierte Historie eher ein Gesprächsarchiv als episodisches
Gedächtnis: Es gibt weder selektiven Abruf noch Konsolidierung und Reflexion.
Generative Agents zeigen ein Gegenmodell, in dem Erfahrungen gespeichert,
dynamisch abgerufen und zu höherwertigen Reflexionen verdichtet werden, die
Ablation der Autoren weist Wahrnehmung, Planung und Reflexion jeweils als
relevante Beiträge aus \parencite{parkEtAl2023}. Eine Übertragung muss dennoch
domänenspezifisch evaluiert werden, weil glaubwürdiges Simulationsverhalten
nicht mit Lernerfolg gleichzusetzen ist.

Zweitens ist Kontrolle überwiegend vor der Aktion (Tool-Whitelist) oder nach
der Antwort (Reviewer) platziert, nicht zwischen Beobachtung und irreversibler
Wirkung. Das System speichert auch kein überprüftes Outcome, aus dem spätere
Entscheidungen lernen könnten. Reflexion nutzt verbales Feedback in einem
episodischen Puffer, um Folgeversuche zu verbessern
\parencite{shinnEtAl2023reflexion}, zugleich hängt dieses Verfahren von der
Qualität der Selbstbewertung ab und garantiert keinen Erfolg. Für StudyMummy
folgt daraus nicht, ungeprüfte Selbstkritik zu akkumulieren, sondern nur
validiertes Reviewer- und Werkzeugfeedback strukturiert zurückzuführen.

\subsection{Verbesserungsvorschlag 1: Begrenztes, konsolidiertes Episodengedächtnis}
\label{sec:verbesserung-eins}

Technisch sollte eine Tabelle \texttt{agent\_episodes} ergänzt werden, deren
Datensatz \texttt{user\_id}, \texttt{session\_id}, \texttt{task\_id}, den
validierten Plan, ausgeführte Aktion, Werkzeugstatus, Reviewer-Urteil, eine
knappe Ergebniszusammenfassung, Zeit und optional ein Embedding enthält. Ein
deterministischer \texttt{MemoryWriter} schreibt erst nach Review und Commit,
interne Gedankengänge werden nicht gespeichert. Ein
\texttt{Memory\allowbreak{}Retriever}
lädt vor der Planung höchstens drei nutzer- und aufgabengefilterte Episoden nach
Semantik, Aktualität und Ergebnisstatus. Parallel fasst ein
\texttt{session\_summary} ältere Chatturns ab einem festen Tokenbudget zusammen,
die letzten Turns bleiben wörtlich. Lernprofil und relevante Confidence Scores
werden als eigene, vertrauenswürdige Felder in \texttt{AgentContext} gebunden.

Dies überträgt CoALAs Trennung von Arbeits-, episodischem und semantischem
Gedächtnis sowie den Abruf früherer Episoden in die vorhandene SQLAlchemy- und
pgvector-Architektur \parencite{sumersEtAl2024}. Erwartet werden stabilere lange
Sitzungen, gezieltere Wiederaufnahme früherer Fehler und begrenzte Promptgröße.
Voraussetzungen sind ein Lösch- und Aufbewahrungskonzept, Migrationen und ein
Evaluationssatz mit Mehrsitzungsverläufen. Zusätzliche Embeddings und Abrufe
erhöhen Latenz und Kosten, Schema, Zusammenfasser und Retrieval-Scoring erhöhen
Komplexität und Wartungsaufwand. Die Sicherheitsoberfläche wächst durch
\emph{Memory Poisoning}: manipulierte Inhalte könnten dauerhaft wiederkehren.
Daher sind Provenienz, Nutzerisolation, TTL, Löschbarkeit und das Speichern nur
validierter Outcomes zwingend. Auch Reflexionen dürfen nicht ungeprüft als
Fakten behandelt werden.

\subsection{Verbesserungsvorschlag 2: Typisierte Plan--Act--Observe--Replan-Schleife}
\label{sec:verbesserung-zwei}

Der Orchestrator sollte ein explizites Objekt \texttt{AgentState} führen. Es
enthält \texttt{plan}, \texttt{observations}, \texttt{attempt},
\texttt{pending\_\allowbreak side\_\allowbreak effect}, \texttt{review} und
\texttt{status}. Der Planner ergänzt pro Aktion eine erwartete Beobachtung und
eine Fehlerstrategie. Zunächst werden nur lesende Werkzeuge ausgeführt, deren
Resultate werden gegen Pydantic-Schemas validiert. Bei Toolfehler oder
Reviewer-Ablehnung erhält der Planner die kompakte Beobachtung und darf höchstens
einmal neu planen. Schreibende Tools werden als \texttt{pending} materialisiert
und erst nach Policy-Prüfung sowie -- beim Kalender -- Nutzerbestätigung
committet. Ein Datenbank-Checkpoint vor der Wirkung ermöglicht sauberen Abbruch,
harte Grenzen für Schritte, Zeit und Toolbudget verhindern Schleifen.

Die Ausgestaltung setzt ReActs alternierende Kopplung von Planen und
Umweltbeobachtung um, ohne Chain-of-Thought offenzulegen
\parencite{yaoEtAl2023react}. Validiertes Reviewer-Feedback kann zusätzlich als
sprachliches Fehlersignal für den nächsten Versuch dienen, entsprechend der
Grundidee von Reflexion \parencite{shinnEtAl2023reflexion}. Der erwartete Nutzen
liegt in kontrollierter Fehlerbehandlung, überprüfbaren Abbruchbedingungen und
dem Schließen der heutigen Lücke, in der Review erst nach Seiteneffekten
erfolgt. Dem stehen mindestens ein weiterer LLM-Aufruf im Fehlerfall, höhere
Latenz und Kosten, komplexere Transaktionen und mehr Zustände im Test gegenüber.
Jeder zusätzliche Rückkopplungskanal kann zudem Injection weitertragen, nur
strukturierte Beobachtungen und gleichbleibende Capability-Grenzen dürfen in
die Neuplanung eingehen.

\subsection{Vergleichende Trade-off-Diskussion}

Das Episodengedächtnis verbessert vor allem Langzeitpersonalisation und
Promptökonomie, vergrößert aber Datenschutz- und Poisoning-Risiken dauerhaft.
Die Replan-Schleife erhöht die Kosten pro problematischem Turn, begrenzt dafür
unmittelbar Fehlwirkungen und macht Toolfehler beobachtbar. Priorität hat daher
Vorschlag~2: Er beseitigt die sicherheitsrelevante Post-hoc-Kontrolle und schafft
den typisierten Outcome, den Vorschlag~1 später speichern kann. Danach sollte
das Gedächtnis zunächst ohne automatisch generierte Reflexionen und nur für
validierte Aufgabenereignisse eingeführt werden. Beide Erweiterungen sind erst
nach Vergleich gegen den bestehenden Workflow sinnvoll, relevante Kriterien
sind Planerfüllung, Toolpräzision, Revisionsrate, Latenz, Tokenkosten und
mehrsitziger Lernfortschritt. Ohne solche Messungen wäre der zusätzliche
Architekturaufwand nicht begründet.
